\scanpage{6}
par $\rho$, mais pas de $\unsure{\mathfrak S_3|\mathcal S_3}$, car on a p.ex
\[
\inserted{\sigma_0\sigma_1\bigl(=\sigma_1\sigma_\infty=\sigma_\infty\sigma_0\bigr)=\rho}
\]
\[
\deleted{[\sigma_0,\sigma_1]\unsure{\sigma_0^{-1}\sigma_1^{-1}}=(\sigma_0\sigma_1)^2=\rho^2}\qquad\text{dans }\unsure{\mathfrak S_3|\mathcal S_3}
\]
mais dans $\operatorname{Aut}(\pi)$, \deleted{\unsure{on a}}
\[
\sigma_0\sigma_1(l_0)=\sigma_0(l_\infty)=l_1=\rho(l_0)
\]
\[
\begin{aligned}
\sigma_0\sigma_1(l_1)&=\sigma_0(l_\infty^{-1}l_0^{-1})=\sigma_0(l_\infty)^{-1}\sigma_0(l_0)^{-1}\\
&=l_1^{-1}(l_1^{-1}l_\infty^{-1})^{-1}=l_1^{-1}l_\infty l_1
\end{aligned}
\]
\[
\begin{aligned}
\sigma_0\sigma_1(l_\infty)&=(l_1^{-1}l_\infty l_1)^{-1}l_1^{-1}=l_1^{-1}l_\infty^{-1}\deleted{\unsure{l_1}}\\
&=l_1^{-1}(l_\infty^{-1}l_1^{-1})l_1=l_1^{-1}l_0l_1
\end{aligned}
\]
donc $\sigma_0\sigma_1=\operatorname{int}(l_1^{-1})\rho$ dans $\operatorname{Aut}(\pi)$ et de $\unsure{=|\bar m}$
\[
\sigma_1\sigma_\infty=\operatorname{int}(l_\infty^{-1})\rho
\]
\[
\sigma_\infty\sigma_0=\operatorname{int}(l_0^{-1})\rho
\]

On a une représentation de $\unsure{\mathfrak S_3|\mathcal S_3}$ dans le groupe \uline{extérieur} à lacets seulement. Notons qu'on ne peut espérer réaliser cette opération extérieure par une vraie opération (du moins pour le groupe discret) vu que $\unsure{\mathfrak S_3|\mathcal S_3}$ opérant sur $\mathbb P^1_{\mathbb C}\setminus\{0,1,\infty\}$ n'a pas de pt fixe. Le fait qu'on ait pu faire opérer $\mathbb Z/3\mathbb Z$ \inserted{sur $\pi$ lui-$=$} provient du fait que son action sur $\mathbb P^1_{\mathbb C}-\{0,1,\infty\}$ admet un (et $\unsure{=|\bar m}$ deux!) point fixe, \inserted{à savoir $\unsure{-\bar j|-j}$ ($j$ racine primitive troisième de $1$)} qu'on \deleted{\unsure{pourra}} a pris comme origine pour calculer le $\pi_1$.

Signalons aussi qu'on a inclus dans le groupe de Teichmüller \inserted{$\Gamma_{0,3}$ opérant sur $\pi_{0,3}$} la ``conjugaison complexe''
\mspage{5}
